\par\Needspace{7\baselineskip}\bigskip
\begin{center}
\pdfbookmark[2]{Olbers an Gauss, Bremen, 7. März 1816}{pallas-brief-18160307}
\textsc{Olbers} an \textsc{Gauss}, Bremen, 7. März 1816.
\end{center}
\noindent\hspace*{\parindent}\SourceLeaf{433}\{. . . . .\quad Jetzt eile ich nur, Ihnen wegen der Preisaufgabe zu Paris, wenn Sie es etwa noch nicht erfahren haben sollten, das Nähere zu melden. Der Ausspruch über den Preis wegen der Perturbationen der Planeten, namentlich der Pallas, ist wieder bis zum Jahr 1817 vertagt worden. Doch müssen die Schriften vor dem 1. October 1816 eingesandt sein.
\par Die Preisfrage war: »Die Theorie der Planeten, deren Excentricität und Neigung zu gross sind, als dass wir im Stande wären, ihre Störungen nach den schon bekannten Methoden genau zu berechnen.« Die Klasse verlangt keine numerische Anwendung, sondern nur analytische Formeln, aber so eingerichtet, dass ein geschickter Rechner fähig sei, sie mit Sicherheit entweder auf den Planeten Pallas, oder auf einen andern der neu entdeckten oder noch zu entdeckenden Planeten anzuwenden. — Es waren nur zwei Abhandlungen eingelaufen, deren Verfasser aber die ausgesprochene Absicht der Klasse in der Preis-Ankündigung nicht genug berücksichtigt haben. Beide (besonders der eine) haben noch zu mancherlei analytische Entwickelungen vorbeigelassen, die die Mathematiker noch erst machen müssten, um sie in den Stand zu setzen, die Auflösung des Problems, die sie gegeben haben, verstehen und beurtheilen zu können. Sie haben es zu sehr versäumt, sich bis zu dem Standpunkte des Calculators herabzulassen, der nun wünschen sollte, Tafeln für die Pallas oder irgend einen andern Planeten zu bilden. Die Nachträge, die sie zu verschiedenen Zeiten eingeschickt haben, sind weit entfernt, alle diese Schwierigkeiten zu heben. Da die Klasse aus diesen Nachträgen und aus den eingeschickten Noten der anonymen Verfasser ersehen hat, dass sie nicht Zeit hatten, sich in alle die nothwendigen Entwickelungen einzulassen, und zugleich in Erwägung zieht, dass auch vielleicht andere Mathematiker, die die Fähigkeit und Geschicklichkeit besitzen, diesen schwierigen Gegenstand zu behandeln, aus derselben Ursache abgehalten worden sind, als Preisbewerber aufzutreten, so hat sie die Preisaustheilung noch bis Januar 1817 ausgesetzt. — Der Preis ist doppelt, eine goldene Medaille, 6000 Francs werth.\quad\mbox{. . . . .\}}
\par
